% sec-01: outline item A.  Figures 1 (mermaid), 2 (d2), 3 (graphviz).
% DRAFT stage.  Three sized slots hold the figure numbers.
\section{The Novel-Performer Case}

\subsection{Three Clauses, One Applicant}

\draftinstr{Open by naming the three candidate clauses and quoting each in one
sentence. Clause 1 is the individual-scientist language of Chapter I, which the
ten prior applications were written under. Clause 2 is the NOVEL PERFORMERS
section of Chapter II. Clause 3 is the SBIR clause, which appears three times:
in the Chapter I summary of recommendations, in Chapter III on deploying SBIR
and STTR strategically to couple federally-seeded companies with the scientific
enterprise, and in Chapter IV on SBIR opening doors for technician-founded
ventures. Quote each from
\appfile{funding/science-golden-age/chunk-01-front-matter-and-summary.md},
\appfile{funding/science-golden-age/chunk-03-chapter-two-revitalizing-science-and-technology-enterprise.md},
\appfile{funding/science-golden-age/chunk-04-chapter-three-securing-dominance-in-critical-and-emerging-technologies.md}
and
\appfile{funding/science-golden-age/chunk-05-chapter-four-science-and-technology-better-lives-of-all-americans.md}.
Cite \texttt{kratsios2026sciencegoldenage}.}

\figslot{mermaid-type / flowchart}{6.4}{Three clauses of one report, four
eligibility tests, and the only\\clause a San Diego firm with no indirect-cost
base survives. The\\two that fail, fail on size, not on merit or on subject
matter.}

\draftinstr{After Figure 1, set Table 4: the three clauses against the four
eligibility tests T1 to T4, with the ChemicalQDevice answer in the last column.
The four tests and their report basis are given in
\appfile{funding/capitalization-plan/mermaid/fig-01-clause-eligibility-filter.md}.
Use five columns: test, question, report basis, this company, verdict.}

\subsection{Why the Pioneer Language Does Not Fit an Entity}

\draftinstr{Two paragraphs. The first states what changed: the applicant is now
a California corporation and not a person, so a person-based mechanism has no
subject to fund. The second is the honest concession, which must not be
softened: the individual-scientist framing is what made the ten prior
applications distinctive and is what the report actually rewards. Name the cost
and forward-reference \S9. Source:
\appfile{funding/pdac-funding-applications/final-apply/sections/sec-01-golden-age-mandate.tex}.}

\subsection{Why the Novel-Performers Chapter Does Not Fit Either}

\draftinstr{One paragraph with three numbers taken from the chapter itself:
tens of millions of dollars, coordinated teams of ten to a hundred people, and
timelines of half a decade. Set them beside this company's \$1,606,000, 2.6 FTE
and 33 months. The conclusion is that the chapter describes a real gap that this
company does not fill, and that claiming otherwise would be the kind of
overreach the report is written against. Source:
\appfile{funding/science-golden-age/chunk-03-chapter-two-revitalizing-science-and-technology-enterprise.md},
the NOVEL PERFORMERS section.}

\figslot{d2-type / true grid}{7.0}{The report's own seven activities, its four
institutional forms,\\and a fifth column it does not carry. An SBIR firm of 2.6
people\\is well suited to two of the seven and unsuited to two others.}

\draftinstr{After Figure 2, one paragraph reading the fifth column: two well
suited, three partly suited, two less suited, and the observation that a firm
claiming to be well suited to all seven would be claiming to be a university, a
corporate lab, a federal lab and an FRO at once. The row-by-row reasoning is in
\appfile{funding/capitalization-plan/d2/fig-02-institutional-form-grid.d2.md}.}

\subsection{The Clause That Fits, and What It Is Worth}

\draftinstr{Open with the structural fact: this company has no negotiated
indirect-cost rate agreement, no facilities to recover, and no administrative
layer to fund. Then price it. Figure 3 and Table 5 carry the arithmetic from
\appfile{funding/capitalization-plan/graphviz/fig-03-indirect-cost-decomposition.gv.md}.
Cite \texttt{nihsbirseed} for the 40 percent allowance and
\texttt{ecfr2026uniformguidance} for the de minimis rate.}

\figslot{graphviz-type / record nodes}{8.4}{The same 1,396,000 of direct work
priced three ways. A university\\F and A rate at 57 percent adds 741,000, the
full SBIR allowance\\adds 695,000, and the rate this plan claims adds 210,000 in
all.}

\draftinstr{Set Table 5 immediately after Figure 3: seven rows, four columns
(line, route A university, route B full allowance, route C this plan), ending in
the premium row. Every figure must reconcile; the checks are listed in
\appfile{funding/capitalization-plan/graphviz/README.md}. Close the section with
two sentences on what the 33.1 percent premium means to a program officer who
has a fixed appropriation.}
